\documentclass[11pt]{exam}
\usepackage{tikz}
\usetikzlibrary {plotmarks}
\usepackage{multicol}
\usepackage{subfigure,amsmath,amssymb}
\usepackage{systeme}

\newcommand{\class}{MATH 96}
\newcommand{\term}{Spring 2025}
\newcommand{\examnum}{Final Exam} 
\newcommand{\examdate}{5/4/2025}
\newcommand{\timelimit}{120 Minutes}

\parindent 0ex

\begin{document} 

\pagestyle{head}
\firstpageheader{}{}{}
\runningheader{\class}{\examnum\ - Page \thepage\ of \numpages}{\examdate}
\runningheadrule

\begin{flushright}
\begin{tabular}{p{2.1in} r l}
\textbf{\class} & \textbf{First and Last Name:}     & \makebox[2.2in]{\hrulefill}\\
\textbf{\term}  &                    &\\
\textbf{\examnum} &  & \\
 & \textbf{Section:} & \makebox[2.2in]{\hrulefill}\\

\end{tabular}\\
\end{flushright}
\rule[1ex]{\textwidth}{.1pt}

This exam contains \numpages\ pages (including this cover page) and
\numquestions\ problems.  Check to see if any pages are missing.  Enter
your name and your section above.
\vspace{4pt}
You may \emph{not} use textbooks, course notes, cell phones, or calculators during this exam.
\vspace{4pt}
You are required to show your work on each problem on this exam, and the following rules apply:
\vspace{4pt}
\begin{itemize}

\item \textbf{Organize your work in a reasonable, neat, and coherent way.} Work that is jumbled, disorganized, or lacks clear reasoning will receive little credit.  

\item \textbf{Partial credit may be given.} Any answer must be supported by calculations, explanation, and/or algebraic work to receive full credit.  Unsupported answers will not receive full credit, but well-argued, incorrect solutions may earn partial credit.

\item \textbf{If you require more space, use the back of a page.} Clearly indicate when you do this below the problem.
\end{itemize}

Do not write in the table to the right.

\hfill
\begin{minipage}[t]{2.3in}
\vspace{0pt}
\gradetablestretch{2}
\vqword{Problem}
\addpoints
\gradetable[v][pages]
\end{minipage}
%=======================================================================================================================
\newpage
\begin{questions}
%=======================================================================================================================
\addpoints
\question[8] Simplify and reduce each expression to a single fraction.
\noaddpoints
\begin{multicols}{2}
\begin{parts}
\part[4] $\displaystyle \left(\frac{1}{6}-\frac{2}{9}\right) \cdot \left(\frac{1}{5}+\frac{1}{4}\right)$
\vspace{2.25in}
\part[4] $\displaystyle \frac{(-4)^2-(-1)^2}{-4^2-(-1^2-2)}$
\end{parts}
\addpoints
\end{multicols}
\vspace{4in}
%=======================================================================================================================
\addpoints
\question[6] Solve~~$\displaystyle \frac{1}{3}\left\vert 2x - 5 \right\vert = 2$.
\vspace{4in}
%=======================================================================================================================
\question[15] Consider the line with the equation $-2x + 5y -10 = 0$.
\noaddpoints
\begin{parts}
\part[6] Find the slope and the $y$-intercepts of the line above.
\vspace{1.5in}
\part[4] Determine whether the line $\displaystyle y -1 = -\frac{5}{2}(x+2)$ is parallel, perpendicular, or neither to the line above.  Explain your reasoning.
\vspace{1.75in}
\part[5] Graph the lines $-2x + 5y -10 = 0$ and  $\displaystyle y -1 = -\frac{5}{2}(x+2)$. \emph{Clearly label two points on each line.}

\begin{figure}[!ht]
\vspace{2pt}
\centering
\begin{tikzpicture}[scale=.9]
\draw[dotted] (-7.5,-4.5) grid (7.5,4.5); 
\draw[<->] (-7.5,0) -- (7.5,0)node[anchor=south] {\footnotesize{$x$}};
\draw[<->] (0,-4.5)--(0,4.5)node[anchor=west] {\footnotesize{$y$}};
\foreach \x in {-7,-6,-5,-4,-3,-2,-1,1,2,3,4,5,6,7}
    \draw (\x cm, 2pt) -- (\x cm, -2pt) node[anchor=north] {\footnotesize{$\x$}};
\foreach \y in {-4,-3,-2,-1,1,2,3,4}
    \draw (2pt, \y cm) -- (-2pt, \y cm) node[anchor=east] {\footnotesize{$\y$}}; 
\end{tikzpicture}
\end{figure}
\end{parts}
\addpoints
%=======================================================================================================================
\newpage
\question[16] Determine the solutions for the provided inequalities and represent the solution sets using interval notation.

\noaddpoints
\begin{parts}

\part[8] $\displaystyle 6 \leq 3\left\vert x - 4 \right\vert$

\vspace{2.5in}
\part[8] $\displaystyle 1 \geq 3x-2 > -1$


\vspace{2.5in}

\end{parts}
\addpoints

%=======================================================================================================================
\question[10] Perform the indicated operations. \emph{Simplify the expression as much as possible.}
\vspace{-8pt}
\noaddpoints
\begin{parts}
\begin{multicols}{2}
\part[5] $\displaystyle \frac{5}{3x} - \frac{2y - 1}{4xy}$
\vspace{4.2in}
\part[5] $\displaystyle \frac{3}{t + 3} - \frac{1}{t - 2}$
\end{multicols}
\vspace{4.2in}
\end{parts}
\addpoints

%=======================================================================================================================
\question[15] Expand the following expressions.
\noaddpoints
\vspace{2.5pt}
\begin{parts}
\part[7] $-(x - 2)(3x + 5)$
\vspace{2in}

\part[8] $(m - 1)(m^2 + 2m + 4)$
\vspace{2in}
\end{parts}
\addpoints

%=======================================================================================================================
\question[10] Find solutions to the given quadratic equations by any appropriate method.
\noaddpoints
\vspace{2pt}
\begin{parts}
\part[5] $3x^2 - 6x = 9$
\vspace{2in}

\part[5] $-x^2 = -3x - 5$
\vspace{2in}
\end{parts}
\addpoints
%=======================================================================================================================

\question[15] For the function $f(x)=2x^2-4x-3$, determine the following.
\noaddpoints
\vspace{2pt}
\begin{parts}
\part[2] $y$-intercept
\vspace{.5in}
\part[5] The vertex form $f(x)=a(x-h)^2+k$, with $a, h, k$ values.
\vspace{2in}
\part[8] $x$-intercepts
\vspace{1.5in}

\end{parts}
\addpoints

\question[5] Find the restricted values of $\displaystyle \frac{x^2 - 9}{x^2 + 2x - 3}$ and simplify the rational expression as much as possible.
\addpoints
\end{questions}
\end{document}


